\textbf{Phần 1. Trắc nghiệm lựa chọn}
\setcounterpageref{ex}{0}

\begin{ex} %[5D2N10]
Số thập phân gồm có hai phần là
\choice
{\True Phần nguyên và phần thập phân}
{Phần nguyên và phần phân số}
{Phần trăm và phần nghìn}
{Phần chục và phần đơn vị}
\loigiai{
Mỗi số thập phân gồm hai phần: phần nguyên và phần thập phân, chúng được phân cách bởi dấu phẩy.\\
\textbf{Đáp án của câu hỏi là A}
}
\end{ex}

\begin{ex} %[5D2N10]
Số thập phân $9{,}17$ được đọc là
\choice
{Chín phẩy một bảy}
{Chín mươi mốt phẩy bảy}
{Chín trăm mười bảy}
{\True Chín phẩy mười bảy}
\loigiai{
Số $9{,}17$ có phần nguyên là $9$ và phần thập phân là $17$, được đọc là: Chín phẩy mười bảy.\\
\textbf{Đáp án của câu hỏi là D}
}
\end{ex}

\begin{ex} %[5D2H10]
Số thập phân gồm $3$ trăm, $2$ chục, $5$ đơn vị, $4$ phần mười, $3$ phần trăm và $1$ phần nghìn được viết là
\choice
{$325{,}43$}
{\True $325{,}431$}
{$32{,}5431$}
{$352{,}431$}
\loigiai{
Phần nguyên là $325$, phần thập phân gồm hàng phần mười là $4$, phần trăm là $3$, phần nghìn là $1$. Số đó được viết là $325{,}431$.\\
\textbf{Đáp án của câu hỏi là B}
}
\end{ex}

\begin{ex} %[5D2H10]
Trong số thập phân $35{,}624$, chữ số $2$ thuộc hàng nào?
\choice
{Hàng chục}
{Hàng phần mười}
{\True Hàng phần trăm}
{Hàng phần nghìn}
\loigiai{
Trong số $35{,}624$, chữ số $6$ ở hàng phần mười, chữ số $2$ ở hàng phần trăm, chữ số $4$ ở hàng phần nghìn.\\
\textbf{Đáp án của câu hỏi là C}
}
\end{ex}

\begin{ex} %[5D2H10]
Giá trị của chữ số $8$ trong số thập phân $2{,}38$ là
\choice
{$8$}
{\True $\dfrac{8}{100}$}
{$\dfrac{8}{10}$}
{$\dfrac{8}{1\,000}$}
\loigiai{
Chữ số $8$ ở hàng phần trăm nên có giá trị là $\dfrac{8}{100}$ (hay $0{,}08$).\\
\textbf{Đáp án của câu hỏi là B}
}
\end{ex}

\begin{ex} %[5D2H10]
Chuyển phân số thập phân $\dfrac{236}{100}$ thành số thập phân ta được
\choice
{$23{,}6$}
{$0{,}236$}
{$236{,}0$}
{\True $2{,}36$}
\loigiai{
Ta có: $\dfrac{236}{100} = 2\dfrac{36}{100} = 2{,}36$.\\
\textbf{Đáp án của câu hỏi là D}
}
\end{ex}

\begin{ex} %[5D2H10]
Phân số thập phân $\dfrac{39}{1\,000}$ được viết dưới dạng số thập phân là
\choice
{\True $0{,}039$}
{$0{,}39$}
{$0{,}0039$}
{$3{,}9$}
\loigiai{
Ta có: $\dfrac{39}{1\,000} = 0{,}039$.\\
\textbf{Đáp án của câu hỏi là A}
}
\end{ex}

\begin{ex} %[5D2H10]
Số thập phân nào dưới đây có cách đọc là: ``Bốn phẩy không năm''?
\choice
{$4{,}5$}
{$4{,}50$}
{$0{,}45$}
{\True $4{,}05$}
\loigiai{
Số thập phân có phần nguyên là $4$, phần thập phân là $05$ được viết là $4{,}05$ và đọc là bốn phẩy không năm.\\
\textbf{Đáp án của câu hỏi là D}
}
\end{ex}

\begin{ex} %[5D2H10]
Viết số đo $118\text{ cm}$ dưới dạng số thập phân với đơn vị mét ta được
\choice
{$11{,}8\text{ m}$}
{$0{,}118\text{ m}$}
{\True $1{,}18\text{ m}$}
{$118{,}0\text{ m}$}
\loigiai{
Ta có: $118\text{ cm} = \dfrac{118}{100}\text{ m} = 1\dfrac{18}{100}\text{ m} = 1{,}18\text{ m}$.\\
\textbf{Đáp án của câu hỏi là C}
}
\end{ex}

\begin{ex} %[5D2H10]
Viết số đo $148\text{ ml}$ dưới dạng số thập phân với đơn vị lít ta được
\choice
{\True $0{,}148\text{ l}$}
{$1{,}48\text{ l}$}
{$14{,}8\text{ l}$}
{$0{,}0148\text{ l}$}
\loigiai{
Ta có: $148\text{ ml} = \dfrac{148}{1\,000}\text{ l} = 0{,}148\text{ l}$.\\
\textbf{Đáp án của câu hỏi là A}
}
\end{ex}

\begin{ex} %[5D2H10]
Số thập phân thích hợp điền vào dấu hỏi chấm trên tia số $0;\; 0{,}1;\; 0{,}2;\; 0{,}3;\; ?;\; 0{,}5$ là
\choice
{$0{,}35$}
{$0{,}04$}
{\True $0{,}4$}
{$0{,}45$}
\loigiai{
Tia số tăng đều $0{,}1$ đơn vị mỗi vạch. Vạch đứng sau $0{,}3$ là $0{,}4$.\\
\textbf{Đáp án của câu hỏi là C}
}
\end{ex}

\begin{ex} %[5D2V10]
Từ bốn tấm thẻ $7$; $0$; $2$ và dấu phẩy ``$,$'', lập được tất cả bao nhiêu số thập phân có phần nguyên gồm một chữ số và phần thập phân gồm hai chữ số?
\choice
{$4$}
{\True $6$}
{$8$}
{$12$}
\loigiai{
Các số thập phân lập được là:
\begin{itemize}
    \item Phần nguyên là $7$: $7{,}02$ và $7{,}20$.
    \item Phần nguyên là $2$: $2{,}07$ và $2{,}70$.
    \item Phần nguyên là $0$: $0{,}27$ và $0{,}72$.
\end{itemize}
Có tất cả $6$ số thập phân thỏa mãn.\\
\textbf{Đáp án của câu hỏi là B}
}
\end{ex}

\textbf{Phần 2. Trắc nghiệm đúng sai}
\setcounterpageref{ex}{0}

\begin{ex} %[5D2H10]
Cho số thập phân $A = 327{,}106$. Xét tính đúng sai của các khẳng định sau:
\choiceTF
{\True Phần nguyên của số $A$ là $327$.}
{Chữ số $6$ trong số $A$ thuộc hàng phần trăm.}
{\True Chữ số $1$ có giá trị là $\dfrac{1}{10}$.}
{\True Số $A$ được đọc là: Ba trăm hai mươi bảy phẩy một trăm linh sáu.}
\loigiai{
\begin{itemize}
    \item Phần nguyên ở bên trái dấu phẩy là $327$.
    \item Chữ số $6$ ở vị trí thứ ba sau dấu phẩy nên thuộc hàng phần nghìn, có giá trị là $\dfrac{6}{1\,000}$ ($\neq$ hàng phần trăm).
    \item Chữ số $1$ ở hàng phần mười nên có giá trị là $\dfrac{1}{10}$.
    \item Số $A$ đọc là: Ba trăm hai mươi bảy phẩy một trăm linh sáu.
\end{itemize}
\begin{itemize}
    \item a) \textbf{Đúng}.
    \item b) \textbf{Sai}: Chữ số $6$ thuộc hàng phần nghìn.
    \item c) \textbf{Đúng}.
    \item d) \textbf{Đúng}.
\end{itemize}
}
\end{ex}

\begin{ex} %[5D2H10]
Xét tính đúng sai của các khẳng định sau về việc chuyển đổi phân số thập phân thành số thập phân:
\choiceTF
{\True Phân số $\dfrac{149}{10}$ được viết dưới dạng số thập phân là $14{,}9$.}
{Phân số $\dfrac{35}{100}$ được viết dưới dạng số thập phân là $3{,}5$.}
{\True Phân số $\dfrac{1\,675}{1\,000}$ viết dưới dạng số thập phân là $1{,}675$.}
{Phân số $\dfrac{7}{20}$ được chuyển thành số thập phân là $0{,}7$.}
\loigiai{
\begin{itemize}
    \item $\dfrac{149}{10} = 14\dfrac{9}{10} = 14{,}9$.
    \item $\dfrac{35}{100} = 0{,}35 \neq 3{,}5$.
    \item $\dfrac{1\,675}{1\,000} = 1\dfrac{675}{1\,000} = 1{,}675$.
    \item $\dfrac{7}{20} = \dfrac{7 \times 5}{20 \times 5} = \dfrac{35}{100} = 0{,}35 \neq 0{,}7$.
\end{itemize}
\begin{itemize}
    \item a) \textbf{Đúng}.
    \item b) \textbf{Sai}: Kết quả là $0{,}35$.
    \item c) \textbf{Đúng}.
    \item d) \textbf{Sai}: Kết quả là $0{,}35$.
\end{itemize}
}
\end{ex}

\textbf{Phần 3. Trả lời ngắn}
\setcounterpageref{ex}{0}

\begin{ex} %[5D2H10]
Viết phân số thập phân $\dfrac{1\,675}{1\,000}$ dưới dạng số thập phân $a{,}bcd$. Tính tổng các chữ số $a + b + c + d$.
\shortans{19}
\loigiai{
\textbf{Đáp số:} $19$\\
$\dfrac{1\,675}{1\,000} = 1{,}675 \Rightarrow a = 1, b = 6, c = 7, d = 5$. Tổng là: $1 + 6 + 7 + 5 = 19$.
}
\end{ex}

\begin{ex} %[5D2H10]
Chuyển phân số $\dfrac{7}{20}$ thành phân số thập phân có mẫu số là $100$. Hỏi tử số của phân số thập phân đó là bao nhiêu?
\shortans{35}
\loigiai{
\textbf{Đáp số:} $35$\\
$\dfrac{7}{20} = \dfrac{7 \times 5}{20 \times 5} = \dfrac{35}{100}$. Tử số là $35$.
}
\end{ex}

\begin{ex} %[5D2H10]
Chuyển phân số $\dfrac{4}{125}$ thành phân số thập phân có mẫu số là $1\,000$. Tử số của phân số thập phân nhận được bằng bao nhiêu?
\shortans{32}
\loigiai{
\textbf{Đáp số:} $32$\\
$\dfrac{4}{125} = \dfrac{4 \times 8}{125 \times 8} = \dfrac{32}{1\,000}$. Tử số là $32$.
}
\end{ex}

\begin{ex} %[5D2H10]
Cái xà đơn cao $165\text{ cm}$. Chiều cao của xà đơn viết dưới dạng số thập phân theo đơn vị mét là $a{,}bc\text{ m}$. Tính giá trị của $a + b + c$.
\shortans{12}
\loigiai{
\textbf{Đáp số:} $12$\\
$165\text{ cm} = 1{,}65\text{ m} \Rightarrow a = 1, b = 6, c = 5 \Rightarrow a + b + c = 12$.
}
\end{ex}

\begin{ex} %[5D2V10]
Từ bốn tấm thẻ $7$; $0$; $2$ và dấu phẩy ``$,$'', lập số thập phân bé nhất có phần nguyên gồm một chữ số và phần thập phân gồm hai chữ số. Hỏi số thập phân bé nhất đó là số nào? (Nhập dưới dạng số thập phân có dấu phẩy).
\shortans{0,27}
\loigiai{
\textbf{Đáp số:} $0{,}27$\\
Các số lập được là $7{,}02; 7{,}20; 2{,}07; 2{,}70; 0{,}27; 0{,}72$. Số bé nhất là $0{,}27$.
}
\end{ex}

\begin{ex} %[5D2V10]
Từ bốn tấm thẻ $7$; $0$; $2$ và dấu phẩy ``$,$'', lập được tất cả bao nhiêu số thập phân có phần nguyên gồm một chữ số và phần thập phân gồm hai chữ số?
\shortans{6}
\loigiai{
\textbf{Đáp số:} $6$\\
Có $3$ cách chọn chữ số phần nguyên ($7, 2, 0$), với mỗi cách có $2$ cách xếp $2$ chữ số còn lại vào phần thập phân $\Rightarrow 3 \times 2 = 6$ số.
}
\end{ex}
